\chapter{Config, and the exception nothing caught}\label{ch:config}

\section{How the config is read}

\texttt{src/ate/config.py} is the whole of ATE's configuration. Four things a developer has to know
before touching it:

\begin{itemize}
  \item \texttt{\_SCHEMA} is the list of accepted keys, grouped by section, each with a type;
    \texttt{\_DEFAULTS} is the default for every one of them. A key the file does not have takes its
    default; a key that is not in the schema is refused with \texttt{unknown config key '<key>' in
    <file>}. \texttt{format} is required and must be 1.
  \item \texttt{<data>} is the only placeholder. It expands to the platform data directory
    (\texttt{/var/lib/ate} on Linux, \texttt{\%PROGRAMDATA\%\textbackslash ATE} on Windows). Any other
    \texttt{<word>} is a config error naming the file and the key, so a typo never becomes a directory
    literally called \texttt{<whatever>}.
  \item Config values get no environment expansion. The platform default path is decided by
    \texttt{\_platform\_default()}, and \texttt{--config <path>} beats it. Note the trap the file records:
    \texttt{sys.platform} is \texttt{win32}, never \texttt{windows}, and \texttt{\_is\_windows} is the only
    platform test used anywhere.
  \item \texttt{ATE\_DEFAULT\_CONFIG} and \texttt{ATE\_DATA\_DIR} are test seams. They are read nowhere
    else in ATE: they exist so a test can exercise ``there is no config at the default location'' in a
    real subprocess without writing to \texttt{/etc/ate} (packet 30c). Do not use them in production code.
\end{itemize}

Adding a key means adding it to \texttt{\_SCHEMA}, to \texttt{\_DEFAULTS}, to the accessor on
\texttt{Config}, and to both manuals --- \texttt{config\_keys()} in \texttt{tests/test\_manuals.py}
derives the set from \texttt{\_SCHEMA} and fails a manual that does not name it.

\section{The application log}

\texttt{src/ate/log.py} opens one \texttt{AppLog} per launch, at
\texttt{<data>/logs/ate-app-<launch time>.log}, honouring \texttt{log.level}. A session given a config
opens its own if the launch did not hand one in; a log handed in belongs to the launch and the session
does not close it. A log directory that cannot be opened must not stop the application: the failure leaves
\texttt{on\_log} as \texttt{None}, and the session falls back to printing.

\section{\texttt{UnhandledExceptionLog}, and what it does not see}\label{sec:exception}

PySide6 does not end the process when a slot, or a \texttt{QThread.run} override, raises. It hands the
exception to \texttt{sys.excepthook} and carries on --- so a worker keeps recording while a screen quietly
stops updating. \texttt{launch()} installs \texttt{UnhandledExceptionLog} on \texttt{sys.excepthook}: it
writes the traceback to the per-launch application log at ERROR, passes the exception to the previous hook
(so stderr and pytest-qt's capture still see it), and emits \texttt{reported} with the log path so the
window can show one line in its status bar:

\begin{lstlisting}[language=ATEShell]
Internal error: this window may have stopped updating. A Test Run in progress keeps recording. Details: <log path>
\end{lstlisting}

There is no dialog, because a modal box would sit in front of a run and nobody has to answer it for the
run to continue, and no abort, because the worker's record is still good. If the log write itself fails,
the Operator still gets the line, with no path claimed.

Know its limits, because they are narrow and were found by test:

\begin{itemize}
  \item it sees a raising slot (direct or queued), a timer callback, and an exception leaving a
    \texttt{QThread.run} override;
  \item it does not see a \texttt{threading.Thread} exception --- those go to \texttt{threading.excepthook},
    and the Device readers, the only such threads, catch their own;
  \item the Test Run path does not use it: \texttt{RunEngine.run}'s guard catches a worker's exception
    first and reports it through the run's own log and Test Report;
  \item \texttt{qInstallMessageHandler} was not used, because it sees Qt's messages, not Python
    exceptions;
  \item it covers the GUI launch only. \texttt{ate \_selftest} and other headless paths have no
    application log and still print.
\end{itemize}

The consequence for whoever writes a test: a child-process test can no longer trust its exit code to see
a slot exception. \texttt{assert\_no\_swallowed\_exception(stderr)} in \texttt{tests/\_fixtures.py} reads
stderr as well, and it is used by the tests that start a Qt process: the children in
\texttt{tests/test\_windows\_ci\_fixes.py}, the two in
\texttt{tests/test\_modals\_never\_block\_headless.py}, the one in
\texttt{tests/test\_gui\_launch\_30a.py}, and the packet 10a tests.

\cando{say where an unhandled exception lands, and write a test that cannot pass by swallowing one.}
